% 张三慧大学物理学第四版第二卷 第28章《原子中的电子》
% 原教材PDF物理页423--426（原书页406--409），配套S PDF373--383。
% 章末正式习题28.1--28.34，勿将思考题28.1--28.14混入。
\section*{第28章\quad 原子中的电子}
\subsection*{章末习题28.1--28.34}
\paragraph{习题 28.1}
求氢原子Balmer（巴耳末）线系的最短与最长波长。
\paragraph{习题 28.2}
氢原子从 \(n=5\) 激发态跃迁至基态，求发射光子的波长，以及原子因反冲所获速率。
\paragraph{习题 28.3}
证明氢原子Bohr能级可等价写成
\(E_n=-\hbar^2/(2m_ea_0^2n^2)\)
以及 \(E_n=-e^2/(8\pi\epsilon_0a_0n^2)\)。
\paragraph{习题 28.4}
证明氢原子 \(n=1\) 的Bohr轨道能量等于
经典质量为 \(m_e\)、电荷为 \(-e\) 的粒子
沿半径 \(a_0\) 圆轨道运行的总能量。
\paragraph{习题 28.5}
由氢原子能级公式推导Balmer线系
\(\widetilde\nu=R(1/4-1/n^2)\)
的Rydberg常量表达式，并求数值，与现代
\(R=1.0973731534\times10^7\,\mathrm{m^{-1}}\) 比较。
\paragraph{习题 28.6}
把正电子和电子组成的正电子偶素视为Bohr原子。
若系统绕质心旋转，总角动量量子化
\(L_n=n\hbar\)，求最小圆轨道半径，
及从 \(n=2\) 到 \(n=1\) 发射光子频率。
\paragraph{习题 28.7}
把地球绕太阳轨道当作Bohr量子化的引力两体系统，
用太阳引力代替库仑力，地球轨道角动量满足
\(L_n=n\hbar\)，运动采用经典近似。
(1) 求圆轨道半径公式；(2) 地球现实轨道
\(r=1.50\times10^{11}\,\mathrm m\) 对应的极大量子数；
(3) 与下一较大可能轨道的半径差。
\paragraph{习题 28.8}
室女座外的一个星系测得两条氢谱线波长
411.7nm与435.7nm，比地球静止氢谱更长。
(1) 判断各对应的氢Balmer跃迁及相应原谱线波长；
(2) 由多普勒红移求星系的退行速率。
\paragraph{习题 28.9}
氢原子的激发态典型寿命
\(\Delta t=10^{-8}\,\mathrm s\)。
利用能量—时间不确定关系估计跃迁回基态时
发射谱线的\emph{自然宽度}，说明线宽定义的系数。
\paragraph{习题 28.10}
氢放电管温度300K时 \(H_\alpha\) 谱线中心频率
\(4.56\times10^{14}\,\mathrm{Hz}\)。
按一维热运动/多普勒模型估计该谱线的频率宽度，
明确RMS宽度、全宽和FWHM定义差别。
\paragraph{习题 28.11}
证明氢原子基态电子的径向概率密度
\(P_r(r)\) 在 \(r=0\) 至无穷的积分为1，
并解释这一归一化的物理意义。
\paragraph{习题 28.12}
求氢原子基态电子离原子核的平均距离
\(\langle r\rangle\)。
\paragraph{习题 28.13}
氢原子处于基态时，求电子库仑势能
\(\langle U\rangle\)，据此计算其平均动能。
并与经典Bohr基态结果比较。
\paragraph{习题 28.14}
氢原子 \(n=2,l=1\)，三个磁量子数
\(m_l=0,+1,-1\) 的 \(2p\) 波函数按教材式给出，
即归一化的
\(\psi_{210}\propto(r/a_0)e^{-r/(2a_0)}\cos\theta\)、
\(\psi_{21\pm1}\propto(r/a_0)e^{-r/(2a_0)}
\sin\theta\,e^{\pm i\varphi}\)。
(1) 求三个定态的概率密度并校对图28.5(b,c)；
(2) 证明三者概率密度之和为球对称；
(3) 验证 \(\psi_{210}\) 对全空间归一化。
\paragraph{习题 28.15}
对轨道量子数 \(l=1\)，求电子轨道角动量矢量
与 \(z\) 轴的允许夹角。
\paragraph{习题 28.16}
氢原子放电时 \(2P_{3/2}\) 与 \(2P_{1/2}\)
的能级劈裂为 \(4.5\times10^{-5}\,\mathrm{eV}\)。
(1) 求Balmer系最低频跃迁线双线的频率差；
(2) 在 \(n=2,l=1\) 电子轨道模型中，
估算电子感受的等效磁场，说明磁矩耦合的约定。
\paragraph{习题 28.17}
在均匀外磁场中求银原子角动量
与外磁场方向可能的夹角及磁势能，
取 \(B=1.2\,\mathrm T\)，价电子近似为
\(s=\frac12,l=0\)。
\paragraph{习题 28.18}
Stern--Gerlach实验中磁极长度4.0cm，
垂直梯度 \(dB/dz=1.5\,\mathrm{T/mm}\)，
银原子炉温2500K。求
(1) 原子所受磁力；
(2) 由磁极流出到记录板形成的
两条银迹间距，按教材给定几何和热速度约定；
忽略碰撞及磁场外的额外漂移。
\paragraph{习题 28.19}
在1.60T磁场中，将核自旋磁矩的反平行/平行
取向以共振交变电磁场翻转。核磁矩沿磁场的分量
大小为 \(1.41\times10^{-26}\,\mathrm{J/T}\)。
求共振频率及电磁波长。
\paragraph{习题 28.20}
应用电子轨道磁量子数与自旋量子数，
证明 (1) 给定 \(n,l\) 的次壳层最多容纳
\(2(2l+1)\) 个电子；(2) 主量子数相同的壳层
最多容纳 \(2n^2\) 个电子。
\paragraph{习题 28.21}
写出基态硼(B,Z=5)、氩(Ar,Z=18)、
铜(Cu,Z=29)、溴(Br,Z=35)的电子排布式。
\paragraph{习题 28.22}
30keV电子轰击靶材产生的X射线截短波长
\(0.041\,\mathrm{nm}\)。用短波限反推普朗克常数。
\paragraph{习题 28.23}
欲产生 \(\lambda=0.100\,\mathrm{nm}\)
的X射线，X光管最低加速电压是多少？
\paragraph{习题 28.24}
40keV电子进入靶材后经历4次碰撞才停止，
假设前三次每次都损失碰前动能的一半。
若每次损失的能量都变为单个X射线光子，
求各次光子的波长。
\paragraph{习题 28.25}
某元素 \(K_\alpha\) X射线波长
\(\lambda=3.16\,\mathrm{nm}\)。
(1) 求K层与L层间能量差；
(2) 依据屏蔽近似Moseley公式估计原子序数。
\paragraph{习题 28.26}
铜K、L层电子离子化能分别为8.979keV
及0.951keV。铜靶特征X射线射到NaCl晶面，
在\emph{掠射角}74.1°处产生一级Bragg极大。
求与晶面平行的相邻Na离子平面间距。
\paragraph{习题 28.27}
CO$_2$ 激光器的光波长10.6\,$\mu$m。
(1) 求CO$_2$分子的这两个振动态能级差；
(2) 在300K热平衡下两能级粒子数比；
(3) 工作时上能级粒子数比下能级多1\%，
按玻尔兹曼形式估算对应形式温度并解释负温度。
\paragraph{习题 28.28}
脉冲激光光脉冲长度10fs，波长500nm。
估算一个脉冲内光振荡的周期数。
\paragraph{习题 28.29}
红宝石脉冲激光器波长694.4nm，脉宽12ps，
单脉冲能量0.150J。
求(1)脉冲空间长度，(2)平均脉冲功率，
(3)脉冲中光子总数。
\paragraph{习题 28.30}
GaAlAs半导体激光器微腔体积200\,$\mu$m$^3$，
连续输出功率5.0mW、波长0.80\,$\mu$m。
求每秒发射多少光子。
\paragraph{习题 28.31}
氩离子激光器光束直径3.00mm，功率5.00W，
波长515nm；经过焦距3.50cm凸透镜
聚焦到毛玻璃上，形成圆形衍射中心亮斑。
(1) 求入射光束平均强度；
(2) 求Airy中心亮斑半径；
(3) 若中心斑集中84\%的总功率，
求中心斑的平均强度。
\paragraph{习题 28.32}
某分子纯转动光谱相邻两谱线的频率差
\(8.6\times10^{10}\,\mathrm{Hz}\)；
已知原子的质量为
\(2.66\times10^{-26}\,\mathrm{kg}\)。
将其近似为对称双原子刚性转子，
求两原子间距离。
\paragraph{习题 28.33}
把氢原子类比为电子围绕固定质子的刚性转子，
转动半径取Bohr半径。
估计从转动基态到第一激发态的能量，
并判断此转动模式对室温气体比热有无显著贡献。
\paragraph{习题 28.34}
CO分子的振动频率为
\(\nu=6.42\times10^{13}\,\mathrm{Hz}\)。
利用两原子约化质量近似求等效振动力常数。
